\documentclass[10pt,a4paper]{amsart}
\usepackage[margin=19mm]{geometry}
\usepackage{amsmath,amssymb,mathtools}
\usepackage[scr=rsfs,cal=euler]{mathalfa}
\usepackage{xfrac}
\usepackage[unicode,hidelinks,bookmarks=false]{hyperref}
\newcommand{\ie}{i.e.\ }
\newcommand{\cf}{cf.\ }
\newcommand{\resp}{respectively\ }
\newcommand{\Subsection}{\subsection}
\providecommand{\nobreakdash}{\nobreak\hbox{-}}
\setlength{\emergencystretch}{2em}
\multlinegap0pt
\usepackage{amssymb}
\usepackage{braket}
\usepackage{bm}
\usepackage{xcolor}
\usepackage{enumerate}
\newtheorem{theorem}{Theorem}
\newcommand{\mbb}[1]{\mathbb{#1}}
\newcommand{\mc}[1]{\mathcal{#1}}
\title{An introduction to monitored quantum systems and quantum trajectories: spectrum, typicality, and phases}
\author{Ryusuke Hamazaki, Ken Mochizuki, Hisanori Oshima, Yohei Fuji}
\date{Source: arXiv:2512.19922v2 (CC BY 4.0)}
\begin{document}
\maketitle

\noindent\emph{Source and licence.} An introduction to monitored quantum systems and quantum trajectories: spectrum, typicality, and phases. Version arXiv:2512.19922v2 (CC BY 4.0), distributed by arXiv under the Creative Commons Attribution 4.0 International licence, http://creativecommons.org/licenses/by/4.0/, as recorded in the arXiv licence metadata for this exact version and shown on the abstract page https://arxiv.org/abs/2512.19922v2. Full licence notice: \texttt{licenses/arxiv-2512.19922-CC-BY-4.0.txt}.\par\smallskip

\noindent\textit{Editorial scope.} Counted unit: the article's theorem theorem (source0.tex lines 4873--4875). The article's complete original proof of it is delivered with it (source0.tex lines 4877--4909, one \texttt{proof} environment of 33 lines). No mathematics is written, repaired or re-derived: the statement and the proof are the article's own lines, in the article's own notation, and no retained line is substituted or re-flowed.  No further source-local context is needed: every cross-reference of every retained line resolves inside the two delivered spans.  Nothing was omitted from any retained span, so the package carries no omission record.  No retained line contains a citation, so the wrapper carries no bibliography and no bibliography entry.  No retained line contains a figure environment, an image-inclusion command or drawing code, so no image is delivered and the package declares no asset dependency.  Editorial formatting only, in addition: an AMS \texttt{amsart} preamble that restates the article's own declarations and macros used by the retained lines, namely the environments \texttt{theorem} on separate counters, together with the standard packages those lines need.  The retained environments print in this document as Theorem 1 in order of appearance, whereas the article prints the same items with the numbering of its own full document.  The complete native source of the article as distributed in the arXiv e-print for this exact version is bundled unchanged as \texttt{sources/191521/source0.tex}.
\begin{theorem}
If $\mc{T}$ is a positive unital linear map, then its operator norm satisfies $\lVert \mc{T} \rVert = 1$.
\end{theorem}

\begin{proof}
We first observe that the operator norm $\lVert \mc{T} \rVert$ is equivalently written as 
\begin{align}
\lVert \mc{T} \rVert = \sup_{\lVert \hat{X} \rVert \leq 1} \lVert \mc{T}[\hat{X}] \rVert.
\end{align}
Suppose that $\hat{X} \in \mbb{B}[\mc{H}]$ has $\lVert \hat{X} \rVert \leq 1$.
Let $\hat{X} = \hat{U}' \hat{\Lambda} \hat{V}'$ be the singular value decomposition of $\hat{X}$, where $\hat{U}'$ and $\hat{V}'$ are unitary matrices and $\hat{\Lambda} = \mathrm{diag}(\Lambda_1, \ldots, \Lambda_d)$ is a diagonal matrix with singular values $\Lambda_1 \geq \cdots \geq \Lambda_d \geq 0$.
Since $\Lambda_1 \leq 1$, we can write $\hat{\Lambda}$ as 
\begin{align}
\hat{\Lambda} 
= \begin{pmatrix} \Lambda_1 && \\ & \ddots & \\ && \Lambda_d \end{pmatrix}
= \frac{1}{2} \begin{pmatrix} e^{i\theta_1} && \\ & \ddots & \\ && e^{i\theta_d} \end{pmatrix} + \frac{1}{2} \begin{pmatrix} e^{-i\theta_1} && \\ & \ddots & \\ && e^{-i\theta_d} \end{pmatrix}
\end{align}
with some $\theta_j \in \mbb{R}$. 
Thus, $\hat{X}$ can be written as $\hat{X} = \frac{1}{2}(\hat{U}_+ + \hat{U}_-)$ with unitary matrices $\hat{U}_+$ and $\hat{U}_-$.
We then find
\begin{align}
\lVert \mc{T}[\hat{X}] \rVert 
= \frac{1}{2} \lVert \mc{T}[\hat{U}_+ + \hat{U}_-] \rVert 
\leq \frac{1}{2} \left( \lVert \mc{T}[\hat{U}_+] \rVert + \lVert \mc{T} [\hat{U}_-] \rVert \right).
\end{align}
We now wish to show $\lVert \mc{T}[\hat{U}] \rVert \leq 1$ for unitary $\hat{U}$.
Using the spectral decomposition $\hat{U} = \sum_{j=1}^d e^{i\phi_j} \hat{P}_j$ with $\phi_j \in \mbb{R}$ and rank-one projectors $\hat{P}_j$, we find
\begin{align}
\begin{pmatrix} \hat{\mbb{I}} & \mc{T}[\hat{U}] \\ \mc{T}[\hat{U}]^\dagger & \hat{\mbb{I}} \end{pmatrix}
= \begin{pmatrix} \hat{\mbb{I}} & \sum_j e^{i\phi_j} \mc{T}[\hat{P}_j] \\ \sum_j e^{-i\phi_j} \mc{T}[\hat{P}_j] & \hat{\mbb{I}} \end{pmatrix}
= \sum_{j=1}^d \begin{pmatrix} 1 & e^{i\phi_j} \\ e^{-i\phi_j} & 1 \end{pmatrix} \otimes \mc{T}[\hat{P}_j],
\end{align}
where we have used $\sum_j \mc{T}[\hat{P}_j] = \hat{\mbb{I}}$ at the last equality. 
Since this matrix is positive semidefinite, by the above lemma we obtain $\lVert \mc{T}[\hat{U}] \rVert \leq 1$. 
This implies $\lVert \mc{T}[\hat{X}] \rVert \leq 1$ for $\lVert \hat{X} \rVert \leq 1$ and thus $\lVert \mc{T} \rVert \leq 1$. 
On the other hand, since $\mc{T}$ is unital, we have $\lVert \mc{T}[\hat{\mbb{I}}] \rVert = \lVert \hat{\mbb{I}} \rVert = 1$ and therefore $\lVert \mc{T} \rVert = 1$.
\end{proof}

\end{document}
